\documentclass[tikz, border=0]{standalone}
\usepackage{fontspec}
\usepackage{unicode-math}
\usepackage{pgfplots}
\pgfplotsset{compat=1.18}

\setmainfont{TeX Gyre Pagella}
\setmathfont{TeX Gyre Pagella Math}
\usepackage{luacode} % Pacote necessário para rodar Lua
\usetikzlibrary{calc,arrows.meta,patterns}

\definecolor{Red}{RGB}{ 159, 36, 48 }
\definecolor{Blue}{RGB}{ 103, 178, 255 }
\definecolor{LBlue}{RGB}{ 186, 220, 255 }
\definecolor{DBlue}{HTML}{132C53} 
\definecolor{gray}{HTML}{708090}
\definecolor{orange}{HTML}{ed9552}
\definecolor{orange2}{HTML}{e29d40}

\pgfmathdeclarefunction{V}{3}{%
        \pgfmathparse{ -( (#2/2)*(abs(#1)^2) + ((#3-#2)/4)*(abs(#1)^4) - (#3/6)*(abs(#1)^6))}%
    }


\tikzset{
    ball style/.style={circle, fill=gray, draw=white, line width=0.15em, inner sep=0.4em},
    noise arrow/.style={<->, >=stealth, thick, gray, shorten <= 2pt, shorten >= 2pt},
    arrowS/.style={->, >=stealth, very thick, line width=0.3em, gray, shorten <= 2pt, shorten >= 2pt},
    jump arrow/.style={->, >=stealth, dashed, very thick, gray, shorten <= 4pt, shorten >= 4pt}
}

% Comando para desenhar a bolinha na posição exata x. 
% Argumentos: #1 = x_pos, #2 = Y_scale, #3 = a, #4 = b, #5 = nome_do_node
\newcommand{\DrawBall}[5]{
    \coordinate (pos) at ({22*(#1)/1.2}, {#2*V(#1,#3,#4)});
    \node[ball style] (#5) at (pos) {};
}
\newcommand{\DrawArrow}[5]{
    \draw[arrowS] ({22*(#1)/1.2}, {#2*V(#1,#3,#4)}) -- ({22*(#1+#5)/1.2}, {#2*V(#1+#5,#3,#4)});
}

% Comando para desenhar as setinhas de flutuação acima da bolinha
\newcommand{\NoiseArrows}[1]{
    \draw[noise arrow] ([yshift=0.8em, xshift=-1.2em]#1.north) to[bend left=10] ([yshift=0.8em, xshift=1.2em]#1.north);
}

\newcommand{\BeginMacros}{
        \def\Lx{22}
        \def\Ly{12}

        \useasboundingbox (-22.5,-14.5) rectangle (22.5,12.5);

        \coordinate (center) at (0,0);

        \coordinate (topC) at (0,\Ly);
        \coordinate (topR) at (\Lx,\Ly);
        \coordinate (topL) at (-\Lx,\Ly);
        \coordinate (topRm) at (0.5*\Lx,\Ly);
        \coordinate (topLm) at (-0.5*\Lx,\Ly);
        
        \coordinate (botC) at (0,-\Ly);
        \coordinate (botR) at (\Lx,-\Ly);
        \coordinate (botL) at (-\Lx,-\Ly);
        \coordinate (botRm) at (0.5*\Lx,-\Ly);
        \coordinate (botLm) at (-0.5*\Lx,-\Ly);

        \coordinate (leftC) at (-\Lx,0);
        \coordinate (leftTm) at (-\Lx,0.5*\Ly);
        \coordinate (leftBm) at (-\Lx,-0.5*\Ly);

        \draw[step=4, dashed, gray] (botL) grid (topR);
}

\newcommand{\EndMacros}{
    \node[white, anchor=north, yshift=-0.5em, font=\bfseries] at (botC) {Magnteização $\symbf{\phi}$};
    \draw[black, thick] (botL) rectangle (topR);
}

\begin{document}

\begin{tikzpicture}[x=1em, y=1em]
    \BeginMacros
    \node[white, anchor=north, yshift=-0.5em, font=\bfseries] at (topC) {Fases do modelo GV};

    \draw[Blue, very thick, line width=0.3em] plot
        [domain=-1.2:1.2, samples=120]
            ({22*\x/1.2}, {50*V(\x,1,-1)});
    \draw[LBlue, very thick, line width=0.3em] plot
        [domain=-1.2:1.2, samples=120]
            ({22*\x/1.2}, {50*V(\x,-1,1)});
    \draw[Red, very thick, line width=0.3em] plot
        [domain=-1.2:1.2, samples=120]
            ({22*\x/1.2}, {120*V(\x,-0.5,2.5)});

    \EndMacros
    
\end{tikzpicture} 

\begin{tikzpicture}[x=1em, y=1em]
    \BeginMacros
    \node[white, anchor=north, yshift=-0.5em, font=\bfseries] at (topC) {Fase Ferromagnética};

    \draw[Blue, very thick, line width=0.3em] plot
        [domain=-1.2:1.2, samples=120]
            ({22*\x/1.2}, {50*V(\x,1,-1)});

    \DrawBall{-0.4}{50}{1}{-1}{ballL}
    \DrawArrow{-0.4}{50}{1}{-1}{-0.1}
    \DrawBall{ 0.4}{50}{1}{-1}{ballR}
    \DrawArrow{0.4}{50}{1}{-1}{0.1}
    \DrawBall{ 0}{50}{1}{-1}{ballC}
    \NoiseArrows{ballC}


    \EndMacros
    
\end{tikzpicture}    
    

\begin{tikzpicture}[x=1em, y=1em]
    \BeginMacros
    \node[white, anchor=north, yshift=-0.5em, font=\bfseries] at (topC) {Fase mista};

    \draw[Red, very thick, line width=0.3em] plot
        [domain=-1.2:1.2, samples=120]
            ({22*\x/1.2}, {120*V(\x,-0.5,2.5)});

    \DrawBall{-0.45}{120}{-0.5}{2.5}{ballL}
    \DrawBall{ 0.45}{120}{-0.5}{2.5}{ballR}
    \NoiseArrows{ballR} 
    \NoiseArrows{ballL} 

    \draw[jump arrow] (ballL.north) to[out=60, in=90] ({0},{120*V(0,-0.5,2.5)});
    \draw[jump arrow] (ballR.north) to[out=60, in=90] ({22/1.2},{120*V(1,-0.5,2.5)});

    \DrawBall{-0.7}{120}{-0.5}{2.5}{ballL}
    \DrawArrow{-0.7}{120}{-0.5}{2.5}{-0.1}
    \DrawBall{ 0.7}{120}{-0.5}{2.5}{ballR}
    \DrawArrow{0.7}{120}{-0.5}{2.5}{0.1}
    
    \EndMacros
    
\end{tikzpicture}


\begin{tikzpicture}[x=1em, y=1em]
    \BeginMacros
    \node[white, anchor=north, yshift=-0.5em, font=\bfseries] at (topC) {Fase Paramagnética};

    \draw[LBlue, very thick, line width=0.3em] plot
        [domain=-1.2:1.2, samples=120]
            ({22*\x/1.2}, {50*V(\x,-1,1)});

    \DrawBall{-0.6}{50}{-1}{1}{0.1}
    \DrawArrow{-0.6}{50}{-1}{1}{0.1}
    \DrawBall{ 0.6}{50}{-1}{1}{0.1}
    \DrawArrow{ 0.6}{50}{-1}{1}{-0.1}
    \DrawBall{ 0}{50}{-1}{1}{ballC}    
    \NoiseArrows{ballC} 
    
    \EndMacros
    
\end{tikzpicture}

\end{document}